\documentclass[a4paper,11pt]{article}

\usepackage[a4paper,margin=2.55cm]{geometry}
\usepackage[T1]{fontenc}
\usepackage[utf8]{inputenc}
\usepackage{lmodern}
\usepackage{microtype}
\usepackage{amsmath,amssymb,amsthm}
\usepackage{booktabs}
\usepackage{longtable}
\usepackage{array}
\usepackage{enumitem}
\usepackage{xcolor}
\usepackage{hyperref}
\usepackage{listings}

\definecolor{LernaBlue}{HTML}{205493}
\definecolor{LernaGreen}{HTML}{2C7A58}
\definecolor{LernaAmber}{HTML}{A86500}
\definecolor{LernaGrey}{HTML}{5E6672}

\newcolumntype{L}[1]{>{\raggedright\arraybackslash}p{#1}}

\hypersetup{
  colorlinks=true,
  linkcolor=LernaBlue,
  citecolor=LernaBlue,
  urlcolor=LernaBlue,
  pdftitle={Lerna Protocol: Claim-UTxO Settlement for Hydra/Cardano eUTxO},
  pdfauthor={Arthur Zhang}
}

\lstdefinestyle{lerna}{
  basicstyle=\ttfamily\small,
  breaklines=true,
  frame=single,
  columns=fullflexible,
  backgroundcolor=\color{gray!4},
  rulecolor=\color{gray!35},
  keepspaces=true
}
\lstset{style=lerna}

\newtheorem{definition}{Definition}
\newtheorem{assumption}{Assumption}
\newtheorem{invariant}{Invariant}
\newtheorem{lemma}{Lemma}
\newtheorem{theorem}{Theorem}
\newtheorem{remark}{Remark}

\newcommand{\Lerna}{\mathsf{Lerna}}
\newcommand{\Claim}{\mathsf{Claim}}
\newcommand{\Head}{\mathsf{Head}}
\newcommand{\State}{\mathsf{State}}
\newcommand{\hash}{\mathsf{H}}
\newcommand{\Value}{\mathsf{Value}}
\newcommand{\Receipt}{\mathsf{Receipt}}
\newcommand{\Proof}{\mathsf{Proof}}
\newcommand{\ada}{\mathsf{ADA}}

\title{\textbf{Lerna Protocol}\\
\large Claim-UTxO Settlement for Hydra/Cardano eUTxO}
\author{Arthur Zhang}
\date{6 June 2026}

\begin{document}
\pagenumbering{Alph}

\begin{titlepage}
\centering
\vspace*{1.1cm}
{\sffamily\bfseries\LARGE Lerna Protocol\par}
\vspace{0.35cm}
{\sffamily\large Claim-UTxO Settlement for Hydra/Cardano eUTxO\par}
\vspace{1.2cm}
{\normalsize Arthur Zhang\par}
\vspace{0.2cm}
{\normalsize 6 June 2026\par}
\vspace{1.2cm}

\begin{abstract}
\noindent
Lerna is an implementation-backed claim-UTxO settlement discipline for
Hydra-family Cardano eUTxO systems. It replaces a prose-only factory exit
story with a concrete claim object, an Aiken spending validator, deterministic
off-chain transition checks, real Hydra/Cardano devnet evidence, and an
executable security-theorem gate. The implemented claim binds each settlement
to one domain, network id, Hydra Head id, factory id, and claim epoch; requires
strict latest-state supersession; conserves lovelace and native assets across
selected settlement outputs and continuation state; emits release receipts for
materialised rights; supports reduced-signature non-interference, native-asset,
and multi-batch profiles; and rejects production acceptance unless local
ledger artifacts hash-match the recorded evidence.

\medskip
\noindent
The claim is intentionally bounded. Lerna is not a mainnet deployment package,
not a replacement for Hydra close/contest/fanout, and not a proof that every
virtual-channel application can settle within arbitrary transaction budgets.
The theorem checked by the repository is conditional on Hydra contestability,
Cardano ledger validation of submitted transactions, proof availability,
watcher/operator liveness, and materialisation capacity inside the committed
claim budget. Under those assumptions, the current implementation makes the
paper's security surface auditable as code: the validator, deterministic model,
real devnet verifier, profile matrix, local file-hash checks, and negative
tests must all cover the same obligations.
\end{abstract}

\vfill
\begin{flushleft}
\textbf{Keywords:} Cardano, Hydra, eUTxO, claim UTxO, Aiken, Plutus, virtual
channels, state channels, non-interference, native assets, release receipts.
\end{flushleft}
\end{titlepage}

\pagenumbering{roman}
\tableofcontents
\clearpage
\pagenumbering{arabic}

\section{Introduction}

Hydra Head gives a fixed participant set a Cardano-like off-chain ledger:
participants deposit UTxOs, sign snapshots, close with a snapshot when
necessary, contest stale closes, and fan out the resolved UTxO set to
L1~\cite{hydra-paper,hydra-docs}. This is already a powerful settlement
substrate. A virtual-channel factory layered inside or near a Head should
therefore avoid inventing an unmeasured parallel dispute game. Its hard problem
is narrower and more operational: when compact rights are no longer being
updated optimistically, what exact ledger object can materialise them without
double release, stale-state settlement, or silent asset loss?

The current Lerna implementation answers with a script-controlled claim UTxO.
A claim UTxO is a compact ledger object carrying unresolved value, latest-state
metadata, authorisation policy, release-receipt state, proof-availability
commitments, and materialisation bounds. Settlement spends the claim through a
compiled Aiken validator. A final settlement consumes all unresolved claim
value. A partial settlement creates a continuation claim UTxO with an inline
datum and a strictly advanced settled-right frontier.

This paper is written to match that implementation. It no longer presents the
claim path as a later research branch, nor does it treat Head-local accounting
alone as the production target. Historical factory-local ideas remain useful as
motivation, but the auditable protocol claim is now the claim-UTxO path checked
by the repository at commit \texttt{66ff116}. The acceptance commands are:

\begin{lstlisting}[language=bash]
npm test
npm run verify:security-theorem
npm run check:devnet
npm run check:profiles
\end{lstlisting}

\paragraph{Contributions.}
The paper and implementation make four concrete contributions.

\begin{description}[leftmargin=2.6cm,style=nextline]
\item[C1. Claim-UTxO state object.]
Lerna defines a claim datum that binds factory identity, Hydra Head identity,
network id, claim epoch, latest-state number, unresolved value, signer policy,
release-receipt root, proof-availability root, and materialisation bounds.

\item[C2. Ledger-enforced settlement transition.]
The Aiken validator checks identity, domain separation, latest-state
supersession, continuation shape, release-receipt progress, authorisation,
output descriptors, selected output indices, exact value conservation, min-ADA,
and continuation/finality conditions.

\item[C3. Profile-backed non-interference and conservation.]
The deterministic model and devnet profile matrix cover reduced-signature
non-interference, native-asset conservation, and multi-batch continuation.
These are not prose-only extensions: each profile has real devnet evidence and
must satisfy the theorem gate.

\item[C4. Executable theorem gate.]
The repository states security obligations as code and fails acceptance unless
the model, validator source, real devnet verifier, local file hashes, profile
matrix, and negative tests jointly cover them.
\end{description}

\section{Layering}

Lerna stays inside the Cardano/Hydra settlement vocabulary. Cardano L1 supplies
eUTxO accounting, validators, inline datums, reference scripts, native assets,
validity intervals, and min-ADA rules~\cite{cardano-eutxo}. Hydra supplies the
outer Head lifecycle: open, snapshot, close, contest, and fanout. Lerna supplies
the inner claim object used when compact virtual-channel rights must be
materialised into ordinary ledger outputs.

This layering creates two design constraints. First, the claim UTxO must not
pretend to replace Hydra contestability. If a stale Head close is not contested
by an actor with the right evidence and authority path, Lerna cannot repair the
outer ledger result by prose. Secondly, compactness must be repaid before final
acceptance. The production report may pass only when the final claim leaves no
unresolved lovelace or native assets.

\begin{remark}[Relation to adjacent work]
Eltoo motivates non-punitive latest-state replacement~\cite{eltoo}. Channel
factories and Perun motivate shared funding and virtual-channel adjudication
patterns~\cite{factories,perun}. Interhead studies virtual Heads across
Heads~\cite{interhead}. Hydrozoa studies a Cardano state-channel design that
separates execution from custody and uses deterministic rollout for large
output sets~\cite{hydrozoa-spec}. Lerna takes these as boundary lessons, not as
drop-in machinery: the current claim is a Cardano/Hydra claim UTxO with
measured evidence.
\end{remark}

\section{System and Threat Model}

\subsection{Actors}

\begin{description}[leftmargin=2.8cm,style=nextline]
\item[Head participant.]
A Hydra party able to run a Hydra node, sign snapshots, close, and contest.

\item[Virtual participant.]
An application endpoint whose right is represented in Lerna claim state. A
virtual participant may or may not be a Head participant.

\item[Operator.]
An application actor that can coordinate batching, submit transactions, or fund
ordinary fees. The operator does not gain authority to redirect claim value.

\item[Watcher.]
An actor that observes Head lifecycle events and Lerna evidence, and can cause
fresh evidence or settlement transactions to be submitted through an authorised
path.
\end{description}

\subsection{Adversary}

The adversary may corrupt participants, operators, watchers, or Head
participants; withhold signatures; provide stale latest-state evidence; propose
same-number conflicting states; withhold proof bodies; try replay across
domains, networks, Heads, factories, or epochs; create outputs below min-ADA;
over-release lovelace or native assets; reorder multi-batch settlement; mutate
continuation datums; or tamper with local evidence files used by production
acceptance.

\subsection{Assumptions}

\begin{assumption}[Hydra contestability]
If the Head is closed with stale state, at least one honest actor can observe
the close and contest with newer valid Head evidence before the contestation
deadline.
\end{assumption}

\begin{assumption}[Cardano ledger validation]
Submitted claim-creation, claim-settlement, and continuation transactions are
validated by the Cardano ledger according to the observed devnet rules.
\end{assumption}

\begin{assumption}[Proof availability]
Proof bodies committed by the proof witness root are available to the verifier
before materialisation is accepted.
\end{assumption}

\begin{assumption}[Watcher and operator liveness]
At least one actor with sufficient evidence and authority can submit accepted
materialisation plans within the relevant Head and ledger deadlines.
\end{assumption}

\begin{assumption}[Bounded materialisation]
The number of outputs, proof bytes, execution units, and continuation value fit
within the budget committed in the claim datum.
\end{assumption}

\section{Claim State}

\begin{definition}[Claim scope]
A claim scope is the tuple
\[
  \Sigma=(\mathsf{domain},\mathsf{networkId},\mathsf{hydraHeadId},
  \mathsf{factoryId},\mathsf{claimEpoch}).
\]
Every latest-state header, transition, receipt, proof-availability commitment,
and non-interference proof is interpreted only inside this scope.
\end{definition}

The implementation uses the following domain separators:

\begin{center}
\begin{tabular}{ll}
\toprule
\textbf{Domain} & \textbf{Purpose}\\
\midrule
\texttt{Lerna:ClaimUTxO:LatestStateHeader} & latest-state and transition binding\\
\texttt{Lerna:ClaimUTxO:ReleaseReceipt} & release receipt nullifiers\\
\texttt{Lerna:ClaimUTxO:ProofAvailability} & proof witness package binding\\
\texttt{Lerna:ClaimUTxO:NonInterferenceProof} & reduced-signature frame proof\\
\bottomrule
\end{tabular}
\end{center}

\begin{definition}[Claim datum]
A claim datum \(C\) contains:
\[
\begin{aligned}
C = (&\mathsf{factoryId},\mathsf{hydraHeadId},\mathsf{claimEpoch},
\mathsf{networkId},\mathsf{domain},\mathsf{latestStateNumber},\\
&\mathsf{stateCommitment},\mathsf{releaseReceiptRoot},
\mathsf{nonInterferenceRoot},\mathsf{authorisationMode},\\
&\mathsf{unresolvedLovelace},\mathsf{unresolvedAssets},
\mathsf{minAdaFloor},\mathsf{requiredSigners},\\
&\mathsf{settledRightCount},\mathsf{lastSettledRightId},
\mathsf{maxOutputCount},\mathsf{maxProofSizeBytes},\\
&\mathsf{maxExecutionMemory},\mathsf{maxExecutionSteps},
\mathsf{proofAvailabilityMode},\mathsf{proofWitnessRoot}).
\end{aligned}
\]
The state commitment is recomputed from length-prefixed canonical fields, not
accepted as an independent assertion.
\end{definition}

The claim value is:
\[
  \Value(C)=
  \mathsf{lovelace}(C.\mathsf{unresolvedLovelace})
  + C.\mathsf{unresolvedAssets}.
\]

\section{Settlement Transition}

A settlement transition spends the current claim UTxO and selects a canonical
batch of rights \(R=(r_1,\ldots,r_k)\) with matching settlement outputs
\(O=(o_1,\ldots,o_k)\). It provides a redeemer \(D\), an authorised signer set
\(S\), and either a final zero-claim result or a continuation claim datum
\(C'\).

\subsection{Admission}

The transition is admitted only if:

\begin{enumerate}[leftmargin=1.2cm]
\item \(D.\mathsf{nextStateNumber} > C.\mathsf{latestStateNumber}\).
\item The domain, network id, Head id, factory id, and claim epoch in the
redeemer match the datum.
\item \(k > 0\) and \(k \le C.\mathsf{maxOutputCount}\).
\item Every selected output is at or above \(C.\mathsf{minAdaFloor}\).
\item Proof size and script-cost evidence are within the committed bounds.
\item The batch starts strictly after \(C.\mathsf{lastSettledRightId}\) in the
canonical settled-right order.
\end{enumerate}

\subsection{Authorisation}

Lerna has two implemented authorisation modes. In full authorisation, the
settlement evidence must contain exactly the datum-bound signer set. In
reduced-signature mode, the authorised signers must be an exact proper subset
and the non-interference witness must show that all materialised outputs belong
to signed owners and that the untouched frame is preserved.

The repository profile matrix currently passes reduced-signature evidence on a
real devnet run. This matters because reduced-signature settlement is a common
source of accidental over-claiming: without a frame proof, an output for an
unsigned owner could be smuggled into a partial settlement.

\subsection{Conservation}

Let \(M\) be the value materialised by the selected settlement outputs and
\(Remainder\) be the value recorded in the continuation claim, or zero for a
final settlement. The validator and deterministic model require:
\[
  \Value(C) = M + Remainder.
\]
This equality is exact over lovelace and native assets. The native-asset
profile is therefore not merely a text extension; it is a profile-matrix run
whose real devnet evidence must pass.

\subsection{Release Receipts}

For each materialised right \(r_i\) and selected output \(o_i\), Lerna emits a
deterministic receipt:
\[
\Receipt_i=\hash(\texttt{Lerna:ClaimUTxO:ReleaseReceipt},
  \Sigma, r_i.\mathsf{rightId}, \mathsf{desc}(o_i), i).
\]
The transition advances the receipt root from the previous root to a new root
covering the settled batch. A later settlement cannot materialise the same
right prefix again because the canonical frontier and receipt root must both
advance.

\subsection{Continuation}

If unresolved value remains, exactly one continuing output must return to the
claim script address with an inline datum \(C'\). That datum must preserve the
scope, signer policy, non-interference root, materialisation budget, proof
availability mode, and proof witness root, while advancing the latest-state
number, receipt root, settled-right count, last settled right id, unresolved
value, and state commitment. If no unresolved value remains, there must be no
continuing claim output.

\section{Validator Enforcement}

The Aiken validator \texttt{validators/lerna\_claim.ak} enforces the settlement
rules at the script boundary. Its key checks are:

\begin{longtable}{L{0.33\linewidth}L{0.57\linewidth}}
\toprule
\textbf{Check family} & \textbf{Validator responsibility}\\
\midrule
\endhead
Input and continuation & Resolve the own input, compare its value to claim
value, require either no continuation for final settlement or exactly one
inline-datum continuation.\\
Identity and domain & Bind datum and redeemer to the same factory, Head,
network, epoch, and latest-state domain.\\
Latest state & Require the old commitment to match the datum and the next state
number to be strictly greater.\\
State commitment & Recompute the current and next state commitments from
canonical fields.\\
Receipt progress & Recompute the expected release-receipt root and require it
to advance.\\
Authorisation & Check exact full signers or exact reduced-signature
non-interference signers.\\
Conservation & Require input claim value to equal selected outputs plus
continuation value across lovelace and native assets.\\
Output descriptors & Bind selected output indices, payment key hashes,
lovelace, and asset bundles to the redeemer descriptors.\\
Admission bounds & Enforce positive numeric bounds, output-count bounds,
min-ADA policy, and continuation-value safety.\\
\bottomrule
\end{longtable}

These checks are intentionally redundant with the deterministic off-chain
model. The model explains and generates expected transitions; the validator
decides whether an on-chain spend is admissible; the real devnet verifier checks
that the claimed local evidence refers to the actual generated artifacts.

\section{Implementation Theorem}

\begin{theorem}[Implementation security theorem]
For any production-accepted claim settlement, if the Hydra/Cardano devnet
evidence, local files, profile matrix, and theorem gate all pass, then the
settlement:
\begin{enumerate}[leftmargin=1.2cm]
\item spends the script-controlled claim UTxO through the compiled claim
validator with matching inline datum, redeemer, and artifacts;
\item binds latest-state evidence to one domain, network id, Head id, factory
id, and claim epoch;
\item strictly supersedes the previous latest-state number;
\item preserves recomputable datum, transition, and continuation commitments;
\item uses exact full authorisation or reduced-signature non-interference;
\item conserves lovelace and native assets across selected outputs and
continuation state;
\item emits deterministic release receipts and advances the settled-right
frontier;
\item respects committed output-count, proof-size, execution-unit, and min-ADA
bounds;
\item binds proof availability before materialisation;
\item accepts only hash-matching local ledger, transaction, artifact,
script-cost, and Hydra lifecycle evidence; and
\item reaches final acceptance only when no unresolved claim value remains.
\end{enumerate}
\end{theorem}

The theorem is not a hand-waved claim. It is encoded in
\texttt{offchain/lerna-security-theorem.mjs} as named obligations. The verifier
fails unless each obligation is covered by the deterministic model, validator
source keywords, strict real devnet checks, local file-hash verification,
profile-matrix evidence, and negative tests.

\begin{center}
\begin{tabular}{ll}
\toprule
\textbf{Obligation id} & \textbf{Meaning}\\
\midrule
\texttt{claim-utxo-ledger-enforcement} & Claim spend uses the compiled validator and matching artifacts\\
\texttt{latest-state-supersession} & Same-scope settlement strictly advances state number\\
\texttt{state-commitment-integrity} & Commitments are recomputable from canonical fields\\
\texttt{exact-authorisation} & Full or reduced signer evidence matches the datum policy\\
\texttt{non-interference-frame} & Reduced settlement affects only signed owners\\
\texttt{asset-conservation} & Lovelace and native assets are exactly conserved\\
\texttt{release-receipt-nullifier} & Materialised rights cannot be claimed twice\\
\texttt{multi-batch-progress} & Partial batches advance a canonical frontier\\
\texttt{materialisation-admission} & Output count, proof size, costs, and min-ADA are bounded\\
\texttt{proof-availability} & Witness availability is bound before materialisation\\
\texttt{ledger-evidence-authenticity} & Local files and ledger observations hash-match the report\\
\texttt{finality-of-accepted-report} & Final acceptance leaves no unresolved claim value\\
\bottomrule
\end{tabular}
\end{center}

\section{Evidence and Evaluation}

The implementation repository contains deterministic evidence, real
Hydra/Cardano devnet evidence, and a profile matrix. Production acceptance uses
four commands:

\begin{lstlisting}[language=bash]
npm test
npm run verify:security-theorem
npm run check:devnet
npm run check:profiles
\end{lstlisting}

\texttt{npm test} runs the Node test suite, including negative tests for stale
state numbers, signer mismatches, conservation failures, receipt mismatches,
output descriptor mutations, and local evidence mismatches.
\texttt{npm run verify:security-theorem} checks the obligation coverage.
\texttt{npm run check:devnet} probes the toolchain, builds the Aiken blueprint,
runs deterministic claim evidence, runs the real Hydra/Cardano devnet flow when
available, refreshes deterministic evidence with the real Head scope, verifies
the real evidence, and asserts acceptance. \texttt{npm run check:profiles}
requires the reduced-signature, native-asset, and multi-batch profiles to pass
real devnet verification.

\begin{center}
\begin{tabular}{llll}
\toprule
\textbf{Profile} & \textbf{Authorisation} & \textbf{Native assets} & \textbf{Batching}\\
\midrule
reduced-signature & reduced non-interference & none & final settlement\\
native-asset & reduced non-interference & exact conservation & final settlement\\
multi-batch & reduced non-interference & none & two batches\\
\bottomrule
\end{tabular}
\end{center}

As of the evidence generated on 6 June 2026, all three profiles report
\texttt{passFailStatus = passed}, \texttt{realDevnetPassFailStatus = passed},
and \texttt{noUnresolvedClaimsAfterSettlement = true}.

\section{Limitations}

The current result is strong enough to change the paper's center of gravity,
but it remains bounded.

\begin{enumerate}[leftmargin=1.2cm]
\item The implementation is local devnet evidence, not a mainnet deployment.
\item The theorem assumes Cardano ledger validation and Hydra lifecycle
observations behave as recorded.
\item Proof bodies committed by the proof witness root must be available to the
verifier.
\item Watcher/operator liveness remains a real deployment assumption.
\item Materialisation is safe only within the committed output-count, proof,
execution-unit, and min-ADA budget.
\item Dynamic membership, arbitrary application state, policy-level mint/burn
logic, and cross-Head settlement composition are outside the accepted claim.
\end{enumerate}

\section{Conclusion}

The implementation changes the right paper. Lerna should now be described by
the object that is actually checked. The auditable object is the claim UTxO: a
script-controlled settlement state
with latest-state binding, exact authorisation or reduced-signature
non-interference, release-receipt nullifiers, exact asset conservation,
multi-batch continuation, materialisation admission, proof availability, and
hash-bound devnet evidence.

The remaining work is not to invent a new roadmap label. It is to keep
tightening the theorem boundary: expand independent review of the Aiken
validator, reproduce devnet evidence on fresh infrastructure, measure mainnet
era limits, and state every deployment assumption at the same precision as the
code checks it.

\appendix
\section{Reproducibility Checklist}

\begin{enumerate}[leftmargin=1.2cm]
\item Check the implementation commit: \texttt{66ff116}.
\item Run \texttt{npm test}.
\item Run \texttt{npm run verify:security-theorem}.
\item Run \texttt{npm run check:devnet}.
\item Run \texttt{npm run check:profiles}.
\item Inspect \texttt{evidence/profile-matrix/profile-matrix.json}.
\item Inspect \texttt{docs/LERNA\_SECURITY\_THEOREM.md} and confirm every
obligation id appears in the theorem gate.
\end{enumerate}

\begin{thebibliography}{99}
\raggedright

\bibitem{hydra-paper}
S. Chakravarty, M. D. Coretti, M. Fitzi, P. Gazi, P. Kant, A. Kiayias, and A.
Russell.
\newblock Hydra: Fast Isomorphic State Channels.
\newblock IACR Cryptology ePrint Archive, Report 2020/299.
\newblock \url{https://eprint.iacr.org/2020/299.pdf}.

\bibitem{hydra-docs}
Cardano Scaling.
\newblock Hydra Head protocol documentation.
\newblock \url{https://hydra.family/head-protocol/docs/protocol-overview}.

\bibitem{hydra-known}
Cardano Scaling.
\newblock Hydra known issues and limitations.
\newblock Consulted 6 June 2026.
\newblock \url{https://hydra.family/head-protocol/docs/known-issues}.

\bibitem{cardano-eutxo}
Cardano Documentation.
\newblock The Extended UTxO accounting model.
\newblock \url{https://docs.cardano.org/about-cardano/learn/eutxo-explainer/}.

\bibitem{interhead}
M. Jourenko, M. Larangeira, and K. Tanaka.
\newblock Interhead Hydra: Two Heads are Better than One.
\newblock IACR Cryptology ePrint Archive, Report 2021/1188.
\newblock \url{https://eprint.iacr.org/2021/1188.pdf}.

\bibitem{hydrozoa-spec}
G. Flerovsky, I. Rodionov, and P. Dragos.
\newblock Hydrozoa: Lightweight state channels for Cardano.
\newblock Technical specification working draft, 7 November 2025.
\newblock \url{https://cardano-hydrozoa.github.io/hydrozoa/hydrozoa.pdf}.

\bibitem{eltoo}
C. Decker, R. Russell, and O. Osuntokun.
\newblock eltoo: A Simple Layer2 Protocol for Bitcoin.
\newblock \url{https://blockstream.com/eltoo.pdf}.

\bibitem{factories}
C. Burchert, C. Decker, and R. Wattenhofer.
\newblock Scalable Funding of Bitcoin Micropayment Channel Networks.
\newblock \url{https://tik-old.ee.ethz.ch/file//49218d6b9325ddb4f18aef8830ec567b/1/scalable_funding.pdf}.

\bibitem{perun}
S. Dziembowski, S. Faust, and K. Hostakova.
\newblock General State Channel Networks.
\newblock ACM CCS 2018.

\bibitem{lerna-impl}
a19q3.
\newblock Lerna claim-UTxO implementation repository.
\newblock Commit \texttt{66ff116}, consulted 6 June 2026.
\newblock \url{https://github.com/a19q3/Lerna}.

\end{thebibliography}

\end{document}
